\scanpage{8}
\mspage{7}

\begin{center}- 55 -\end{center}
\margin{$\mathrm{n}^{\mathrm{o}}357$}

[N.B. -- Ce n'est pas par amour des méthodes éculées et pisseuses que le rédacteur, dans la construction de la mesure de Haar, est revenu aux mesures à support \uline{fini} au lieu des mesures à support \uline{compact} des rédactions précédentes ; c'est en vertu du rasoir d'Occam (``entitates non sunt multiplicandae''), et pour pouvoir laisser reposer en paix dans la Remarque de la fin du Chap.VII, § 2, n°2, la théorie générale de l'opération $\mu*f$, ou $\gamma(\mu)f$, lorsqu'on ne s'est pas donné de jauge ; ce n'est évidemment pas difficile (cf. loc. cit.), mais ça paraît bien superflu $\ldots$]

On écrira $\mathcal{K}$, $\mathcal{K}(K)$, etc. au lieu de $\mathcal{K}(G)$, $\mathcal{K}(G,K)$, etc. ; on notera $\mathcal{K}_+$ l'ensemble des fonctions $f\in\mathcal{K}$ à valeurs $\geq0$, $\mathcal{K}^*_+$ l'ensemble des fonctions $g\in\mathcal{K}_+$ autres que $0$, et $\mathcal{K}^*_+(K)$, si $K$ est une partie compacte de $G$, l'ensemble des fonctions $g\in\mathcal{K}^*_+$ à support contenu dans $K$. Pour $f\in\mathcal{K}_+$, $g\in\mathcal{K}^*_+$, désignons par $A(f,g)$ l'ensemble des mesures positives $\mu\neq0$, à support fini, telles que $f\leq\gamma(\mu)g$. \uline{Cet ensemble n'est pas vide.} En effet, soient $K$ le support de $f$, et $U$ un ouvert non vide tel que $m=\inf_U(g)>0$ ; soit $M=\sup_K(f)$. On peut recouvrir $K$ par un nombre fini d'ensembles $s_iU$ ; alors la mesure $(M/m)\sum\varepsilon_{s_i}$ est dans $A(f,g)$. Posons maintenant :
\[
(f:g)=\inf_{\mu\in A(f,g)}\Vert\mu\Vert.
\]
Il est clair qu'on a $0\leq(f:g)<+\infty$. De plus, ce symbole a les propriétés suivantes :

(i) $(\gamma(\varepsilon_s)f:g)=(f:g)$ quel que soit $s\in G$.

(ii) $(cf:g)=c.(f:g)$ quel que soit $c\geq0$.

(iii) $f\leq f'$ entraîne $(f:g)\leq(f':g)$.

(iv) $(f:g)\geq(\sup f)/(\sup g)$.

(v) $(f:h)\leq(f:g).(g:h)$ pour $g\in\mathcal{K}^*_+$, $h\in\mathcal{K}^*_+$.

(vi) si $f$, $f_0$, $g$ sont dans $\mathcal{K}^*_+$, on a :
